%------------------------------------------------
% Latex-Grundgerüst für Seminarausarbeitungen
%
% zu Erzeugen mit
%   pdflatex ausarbeitung.tex
%
% von Carsten Gutwenger, Nils Kriege, Denis Kurz
% 
%------------------------------------------------

\documentclass[a4paper,12pt,twoside,headsepline]{scrartcl}
\usepackage[utf8]{inputenc}

\usepackage[automark]{scrlayer-scrpage}                     % Kopf-/Fußzeilen
\usepackage[ngerman]{babel}                                 % deutsche Sprache
\usepackage[T1]{fontenc}                                    % Unterstützung für Umlaute mit Fonts

\usepackage[pdftex]{graphicx,color}                         % Einbetten von Grafiken, Farbe
\usepackage{subcaption}                                     % subfigures
\usepackage{url}                                            % \url{}-Kommando
\usepackage[pdftex,pdfpagelabels]{hyperref}                 % Hyperlinks im PDF
\usepackage[round]{natbib}                                  % Literaturverzeichnis
\usepackage{doi}                                            % DOIs

\usepackage[section,boxed]{algorithm}                       % Floating-Umgebung für Algorithmen
\usepackage{algpseudocode}                                  % Pseudocode für Algorithmen
\usepackage{amsmath}                                        % Mathe-Befehle: \frac{m}{n}, \log, \mathcal O, ...
\usepackage{amssymb}                                        % Mathe-Symbole: \mathbb N, ...
\usepackage{amsthm}                                         % Theorem-Umgebungen, siehe \newtheorem und letzte Subsection
\usepackage{tikz}                                           % Zeichnungen, z.B. von Graphen
\usepackage{todonotes}                                      % kaum zu übersehende TODO-Anmerkungen

% Header und Footer
\pagestyle{scrheadings}
\ofoot{}\cfoot{}\ifoot{}
\lehead{\pagemark}
\rehead{\slshape\leftmark}
\lohead{\slshape\leftmark}
\rohead{\pagemark}

\newtheorem{definition}{Definition}
\newtheorem{beispiel}{Beispiel}
\newtheorem{bemerkung}{Bemerkung}

\newtheorem{proposition}{Proposition}
\newtheorem{lemma}{Lemma}
\newtheorem{satz}{Satz}
\newtheorem{theorem}{Theorem}
\newtheorem{korollar}{Korollar}
\renewcommand*\baselinestretch{1.35}\selectfont

\begin{document}

%------------------------------------------------
% Titelseite
%------------------------------------------------

\begin{titlepage}
\sffamily
\vspace*{-3cm}
\includegraphics[width=8cm]{tud_logo_rgb}

\vspace*{4cm}
\begin{center}
\large
{\Large Proseminarausarbeitung} \\[1ex]
{\LARGE\textbf{Operationelle Semantik}} \\[3ex]
Abdulrahman Alrefai \\[1ex]
\today \\[7ex]
im Rahmen des Seminars \\[1ex]
{\Large\textbf{Software Verification}} \\[1ex]
von Prof. Dr. Falk Howar \\[1ex]
Sommersemester 2019
\end{center}


\noindent
\textbf{Basierend auf:}\\
Glynn Winskel. The formal semantics of programming languages: an introduction, \textit{MIT press}


\vspace*{\fill}
\noindent
Fakultät für Informatik\\
Lehrstuhl 14 -- Software Engineering\\
TU Dortmund
\end{titlepage}

%------------------------------------------------
% Inhaltsverzeichnis
%------------------------------------------------

\tableofcontents
\clearpage



%------------------------------------------------
% Einführung
%------------------------------------------------

\section{Einführung}

Die Bedeutung der präzisen Semantik für Programmier- und Spezifikationssprachen wurde seit den sechziger Jahren mit der Entwicklung der ersten höheren Programmiersprachen bekannt. 
Beim Erlernen einer neuen Programmiersprache wird normalerweise nur die Syntax angeschaut, um bestimmte Sprachkonstrukte einsetzen zu können. Syntax behandelt jedoch nur korrekt gebildete Sätze. 
Dies bedeutet, dass der Programmierer überfordert wird, wenn es notwendig wird, zu überprüfen, ob das angegebene Programm tatsächlich den beabsichtigten Vorgang durchläuft.Beim Programmierung kommt es oft vor, dass trotz eines syntaktischen korrekten Programms ein falsches Ergebnis zurückgeliefert wird.  .Der Grund dafür ist, dass das Programm nicht unter allen Umständen semantisch korrekt ist. Programme sind normalerweise zu groß und zu komplex, um sie vollständig zu verstehen und zu überprüfen; daher ist es wichtig, dass man ein tiefes Verständnis für operationale Semantik hat.
 Operationale Semantik erzeugt ein gekennzeichnetes Übergangssystem, dessen Zustände die geschlossenen Begriffe über einer algebraischen Signatur beschreiben und deren Übergänge zwischen Zuständen induktiv aus einer Sammlung von Übergangsregeln der Form $\frac{ Permises} {Conclusion}$ erhalten werden.Wenn alle Prämissen wahr sind, stimmen auch die Schlussfolgerungen

\section{Zustände}
Die Semantik von Programmiersprachen beschäftigt sich mit der Bedeutung von Programmen, die auf einem Computer ausgeführt werden sollen und die verschiedenen verfügbaren Ressourcen nutzen.
 Der Status des Speichers während der Ausführung muss berücksichtigt werden, da der Inhalt des Speichers sich ändern kann; dies geschieht beispielsweise, wenn neue Datenbank angelegt werden oder Variablen aktualisiert werden.
Außerdem kann der Inhalt einer Variable im Moment des Programmstarts  einen großen Einfluss auf das Endergebnis haben.
Nur ein Teil vom Speicher ist von Interesse (Werte der Variablen ).
Der Status kann folgendermaßen definiert werden.Es ist eine Funktion, die Variablen Werte  gibt.
 Die Menge der Zustände $\Sigma$ besteht aus Funktionen $\sigma$ : Loc $\rightarrow$ N.Somit ist $\sigma$ (X) der Wert oder Inhalt der Stelle X im Zustand $\sigma$ .
 
 \subsection{Operationelle Semantik}
Operationelle Semantik ist die Bedeutung eines Programms von einem bestimmten Zustand (Startzustand) und Beobachtung, wie Status im Ende der Ausführung des Programms geändert wird (Endzustand).

\section{Operationelle Semantik für die einfache imperative Sprache (IMP)}
IMP ist ein Programm aus einer Folge von Anweisungen besteht, die vorgeben, in welcher Reihenfolge was vom Computer getan werden soll.
 Verhalten des IMP  wird durch Regeln beschrieben, die festlegen, wie seine Ausdrücke ausgewertet und seine Befehle ausgeführt werden.
In unseren Beispiel beziehen sich die Ausdrücke dabei auf

n $\in$ Num, wobei Num positive und negative ganze Zahlen mit Null ist,\\
b $\in$ BExp, wobei BExp Boolesche Ausdruck ist,\\
a $\in$ AExp, wobei AExp Arithmetic Expression ist,\\
c $\in$ commands und\\
X $\in$ Loc    (Location).

Ein arithmetischer Ausdruck besteht aus einer Folge von arithmetischen Operanden (Konstanten, Variablen), die durch arithmetische Operatoren und Klammern voneinander getrennt oder miteinander verknüpft sind.\\ 
Ein boolescher Ausdruck ist ein Ausdruck, der entweder wahr oder falsch ist.


		
			
\textbf{Formationsregeln}:Formationsregeln sind abstrakte Syntax , indem neue Ausdrücke erstellt werden.Nur die Bedeutung von Programm ist von Interesse, also nicht wie  das aufgeschrieben wird.\\
Hier beziehen die Formationsregeln sich für unsere Ausdrücke dabei auf\\
	
a::= n $\mid$ $a_0$ $\oplus$ $a_1$,\\
b::= true $\mid$ false $\mid$ $a_0$ == $a_1$ $\mid$ $a_0$ $\geq$ $a_1$ $\mid$ $a_0$ $\leq$ $a_1$ $\mid$ $\neg$ b $\mid$ $b_1$ $\vee$ $b_0$ $\mid$ $b_1$ $\wedge$ $b_0$,\\   
c ::= skip $\mid$ X := a $\mid$  $c_0$, $c_1$ $\mid$ if b then $c_0$ else $c_1$ $\mid$ while b do c und \\
$\oplus$ = + $\mid$ - $\mid$ *.	\\
		
Hinweise :
\begin{itemize}
\item wenn $a_0$ und $a_1$ arithmetische Ausdrücke sind, ist es auch die Resultat von $a_0$ $\oplus$ $a_1$ usw. 
\item (\textbf{::=}) sollte gelesen werden als " kann sein " .
		\item (\textbf{$\mid$}) sollte gelesen werden als (oder) .
\end{itemize} 

\subsection{Syntaktische Menge }
Seien $a_0$ , $a_1$ Elemente  $\in$ gleiche syntaktische Menge. Die Notation "$a_0$ == $a_1$"heißt, dass $a_0$ und $a_1$ den gleichen Wert haben oder identisch sind. 
Die  arithmetische Ausdruck "12 - 2"  wird aus ganzen Zahlen aufgebaut. Jedoch sie sind nicht syntaktisch identisch mit 10 , obwohl die Ergebnisse identisch sind.

\section{Auswertung}
Arithmetische Ausdrucke wird in Integer ausgewertet und boolesche Ausdrucke wird in True oder False  ausgewertet. Aber kommt es die Frage, wie ein Programm sich verhaltet?
Erstens wird den Status definiert wird, dann werden boolesche Ausdrucke und arithmetische Ausdrucke ausgewertet. Abschließend wird das Kommando ausgeführt.

\subsection{Auswertung von Arithmetik Expression}
Die Auswertung eines arithmetischen Ausdrucks a in einem Zustand  $\sigma$ to n.

\begin{center}  (a,$\sigma$) $\to$ n.    
\end{center}
                      
Auswertung eines Arithmetik Expression  ($a_0$ $\oplus$ $a_1$):

Zuerst wird das Element $a_0$  mit $n_0$ ausgewertet, dann wird das gleiche Verfahren für das Element $a_1$ angewendet.
Zuletzt wird die Notation "$n_0$ $\oplus$ $n_1$"  berechnet und n wird als Ergebnis von der Notation "$a_0$ $\oplus$ $a_1$" erhaltet.

Also werden die folgende Regeln für die Auswertung benötigt,

\begin{enumerate}

\item Bei der Auswertung von Nummern : wird davon ausgegangen, dass diese bereits ausgewertet sind.\label{1}
\begin{center}
(n,$\sigma$) $\to$ n.    

 und so keine premises braucht   $\frac{ } {\text{(n,$\sigma$) $\to$ n}}$
\end{center}  
\item Bei der Auswertung von Location : wird in seiner Inhalt ausgewertet .\label{2}

\begin{center}
	(X,$\sigma$) $\to$ $\sigma(X)$.    
	
	und so keine premises braucht    $\frac{ } {\text{(X,$\sigma$) $\to$ $\sigma(X)$}}$
\end{center} 

\item \label{3} Bei der Auswertung von Operationen, wobei n die Ergebnis von $n_0$ $\oplus$ $n_1$ ist.

\begin{center}
	 $\frac  {\text {($a_0$ , $\sigma$) $\to$ $n_0$   ($a_1$ , $\sigma$) $\to$ $n_1$} } { \text{($a_0$ $\oplus$ $a_1$,$\sigma$) $\to$ n} }$ 
\end{center}


\end{enumerate}   
Hinweis: Regeln ohne premises : heißt Axiome. 



Beispiel : \text{((Init+30)-(7+9),$\sigma$)$\to$ 14 }
\begin{center} \fontsize{18}{15}
	
	$\frac{  \text{  $\frac { \text {$\frac  {}  {\text{ (Init,$\sigma$)$\to$0 $\ref{2}$ $\quad$          (30,$\sigma$)$\to$ 30 $\ref{1}$ }}$}   } {\text{(Init+30,$\sigma$)$\to$30  $\ref{3}$}}$ $\qquad$   $\frac { \text {$\frac  {}  {\text{ (7,$\sigma$)$\to$7 $ $\ref{1}$\quad$           (9,$\sigma$)$\to$9 $\ref{1}$}}$}} {\text{(7+9,$\sigma$)$\to$16 $\ref{3}$}}$  }} {\text{((Init+30)-(7+9),$\sigma$)$\to$ 14 $\ref{3}$}}$.
	
	
	
	*Hinweis:Init ist eine location mit $\sigma$(Init) = 0. 
\end{center}

Der Stil der operationalen Semantik, die hier genutzt wird, ist strukturell und wird auch big-step semantics genannt.
Eine strukturelle Semantik ist syntaxorientiert und induktiv.
Die Semantik eines Programms wird durch die Semantik seiner Teile definiert.

 strukturell ist Syntaxorientiert und induktiv und  die Semantik eines Programms wird durch  die Semantik seiner Teile definiert.
              
\subsection{Auswertung von Boolean Expression}  

unsere ziel ist die Boolean Expression in true oder false Value auszuwerten . 

     
 Die Regeln bestehen aus, 
 \begin{enumerate}
 \item \label{11} (true,$\sigma$)$\to$true, 
	\item \label{12} (false,$\sigma$)$\to$false ,
	\item \label{13}  $\frac  {\text {($a_0$ , $\sigma$) $\to$ $n_0$   ($a_1$ , $\sigma$) $\to$ $n_1$} } { \text{(($a_0$==$a_1$),$\sigma$) $\to$ true} }$  
	, wenn $a_0$ gleich $a_1$ ist ,
	\item \label{14}  $\frac  {\text {($a_0$ , $\sigma$) $\to$ $n_0$   ($a_1$ , $\sigma$) $\to$ $n_1$} } { \text{(($a_0$==$a_1$),$\sigma$) $\to$ false} }$
	, wenn $a_0$ ungleich $a_1$ ist,
	\item \label{15}  $\frac  {\text {($a_0$ , $\sigma$) $\to$ $n_0$   ($a_1$ , $\sigma$) $\to$ $n_1$} } { \text{(($a_0$ $\geq$ $a_1$),$\sigma$) $\to$ true} }$
	, wenn $a_0$ gleich oder grosser als $a_1$ ist,
	\item \label{16}  $\frac  {\text {($a_0$ , $\sigma$) $\to$ $n_0$   ($a_1$ , $\sigma$) $\to$ $n_1$} } { \text{(($a_0$ $\geq$ $a_1$),$\sigma$) $\to$ false} }$
	, wenn $a_0$ ungleich oder kleiner als $a_1$ ist,
	
	\item \label{17} $\frac  {\text {($b$ , $\sigma$) $\to$ $true$ } } { \text{($\neg$b,$\sigma$) $\to$ flase} }$,
	\item \label{20} $\frac  {\text {($b$ , $\sigma$) $\to$ $false$ } } { \text{($\neg$b,$\sigma$) $\to$ true} }$,
	
	\item \label{18} $\frac  {\text {($b_0$ , $\sigma$) $\to$ $t_0$   ($b_1$ , $\sigma$) $\to$ $t_1$} } { \text{(($b_0$ $\land$ $b_1$),$\sigma$) $\to$ t} }$ und
	\item \label{19} $\frac  {\text {($b_0$ , $\sigma$) $\to$ $t_0$   ($b_1$ , $\sigma$) $\to$ $t_1$} } { \text{(($b_0$ $\vee$ $b_1$),$\sigma$) $\to$ t} }$.	
\end{enumerate}

Ein Beispiel der Anwendung dieser Regeln wird im Folgenden durch \\
\begin{center}
$\frac  {\text {($b_0$ , $\sigma$) $\to$ $true$ $\ref{17}$   ($b_1$ , $\sigma$) $\to$ $false$ $\ref{20}$} } { \text{(($b_0$ $\vee$ $b_1$),$\sigma$) $\to$ true $\ref{18}$} }$ gegeben.\\
\end{center}
Seien $b_0$ , $b_1$ Elemente  $\in$ gleiche syntaktische Menge (BExp).
In unseren Beispiel wird $b_0$ ausgewertet , dann $b_1$ ausgewertet ,abschließend wird die Ergebnis von der Notation "$b_0$ $\vee$ $b_1$"  berechnet . 
Bei der Auswertung einer der beide Elemente mit true  sollte das andere Element nicht auswerten.So  das gesamte Ergebnis ist true, solange $b_0$ true ist. 
Aus diesem Grund sind unsere Auswertung-Regeln nicht optimal. 
Es ist empfehlenswert, dass mehr effizient Regeln benutzt werden.
Ein Beispiel für diese Regeln wird so angegeben :


\begin{center}

$\frac  {\text {($b_0$ , $\sigma$) $\to$ $true$ } } { \text{($b_0$ $\vee$ $b_1$,$\sigma$) $\to$ true} }$,\\
$\frac  {\text {($b_0$ , $\sigma$) $\to$ $false$   ($b_1$ , $\sigma$) $\to$ $true$} } { \text{(($b_0$ $\vee$ $b_1$),$\sigma$) $\to$ true} }$ und\\
$\frac  {\text {($b_0$ , $\sigma$) $\to$ $true$   ($b_1$ , $\sigma$) $\to$ $true$} } { \text{(($b_0$ $\vee$ $b_1$),$\sigma$) $\to$ true} }$.
\end{center}


\section{Ausführung von Kommandos}  
Der Ziel der Ausführung von Kommandos ist den Zustand von Variablen  zu ändern.
(c,$\sigma$)$\to$ $\sigma_0$ : Die Ausführung des Befehls c im Zustand $\sigma$ und  beendet im Endzustand $\sigma_0$.\\
(X:=1,$\sigma$) $\to$ $\sigma$[1/X]: Nach der Ausführung der Zuweisung zu Position X, wird den Wert von $X$ durch 1 ersetzt.


 Hier ist die notwendige \textbf{Regeln}
\begin{enumerate}
	\item (skip,$\sigma$) $\to$ $\sigma$ ,
	\item $\frac  {\text {(a , $\sigma$) $\to$ m} } { \text{ (X:=a,$\sigma$) $\to$ $\sigma$[m/X]} }$,
	\item  $\frac  {\text {($c_0$ , $\sigma$) $\to$ $\sigma_0$   ($c_1$ , $\sigma_0$) $\to$ $\sigma_1$} } { \text{($c_0$;$c_1$,$\sigma$) $\to$ $\sigma_1$} }$,
	\item $\frac  {\text {(b , $\sigma$) $\to$ true   ($c_0$ , $\sigma$) $\to$ $\sigma_1$} } { \text{(if b then $c_0$ else $c_1$,$\sigma$) $\to$ $\sigma_1$} }$,
	\item $\frac  {\text {(b , $\sigma$) $\to$ false   ($c_1$ , $\sigma$) $\to$ $\sigma_1$} } { \text{(if b then $c_0$ else $c_1$,$\sigma$) $\to$ $\sigma_1$} }$,
	\item $\frac {\text{(b,$\sigma$) $\to$ false}} {\text{(while b do c, $\sigma$)$\to$ $\sigma$ }}$ und
	\item $\frac {\text{(b,$\sigma$)$\to$ true $\quad$ (c,$\sigma$)$\to$ $\sigma_2$  $\quad$ (while b do c ,$\sigma_2$) $\to$ $\sigma_1$}}{\text{(while b do c,$\sigma$) $\to$ $\sigma_1$}}$  . 
	              
      \end{enumerate}
  
  \section{Zusammenfassung}
  Es gibt noch viele Semantik Stil Z.B small-Step Semantik , Reduktion semantik .
  so was haben wir hier genutzt ist Big-step-semantik . Außerdem war es nicht immer die beste Lösung Z.B in Boolean Expression . Es hängt davon ab , welche Eigenschaften  gebraucht werden . 
  
  \subsection{Small-step-Semantik vs Big Step Semantik }  
  \textbf{small-step-Semantik} :
  \begin{itemize}
  	\item Hier interessieren wir für jede einzelnen kleinen Berechnungsschritten
  	\item  die einmalige Bewertung ist nützlich , um bestimmte Eigenschaften zu beweisen.
  	\item Reihenfolge der Auswertung .
  	\item Small-Step-Semantik Beispiel :  $\frac  {\text {($a_0$ , $\sigma$) $\to$ $a`_0$ } } { \text{($a_0$+$a_1$,$\sigma$) $\to$ $a`_0$+$a_1$} }$ $\quad$  $\frac  {\text {($a_1$ , $\sigma$) $\to$ $a`_1$ } } { \text{($n_0$+$a_1$,$\sigma$) $\to$ $n_0$+$a`_1$} }$
  \end{itemize}
  
  
  
  \textbf{big-step-Semantik}
\begin{itemize}
	\item Verhalten der Teile eines Programmkonstrukts dessen Gesamtverhalten (Teile und Herrsche).
	\item weinger Regeln $\rightarrow$ schneller im Beweis.
\end{itemize}
  
  

%------------------------------------------------
% Literaturverzeichnis
%------------------------------------------------

\nocite{*} 
\bibliographystyle{abbrvnat-ger}
\bibliography{literatur}
\addcontentsline{toc}{section}{\bibname}


\end{document}
